\chapter{Baselines and Comparative Models}
\label{ch:baselines}

\section{Legacy Multi-LLM Benchmark (Model Selection Phase)}

The \texttt{SkillLLM\_v2} benchmark trained and evaluated five open-weight models on SkillSpan with shared prompts, splits, and evaluation code (\texttt{configs/skillspan\_baselines.yaml}).
Table~\ref{tab:legacy-llm} lists completed test metrics from \texttt{results/skillspan\_baselines/summary\_report.json}.

\begin{table}[htbp]
  \centering
  \caption{Legacy three-epoch SkillSpan benchmark (3{,}569 test examples). Gemma 4 26B did not complete.}
  \label{tab:legacy-llm}
  \begin{tabular}{lrrrr}
    \toprule
    Model & Exact F1 & Token F1 & Partial F1 & Status \\
    \midrule
    Mistral Small 3.1 24B & 0.665 & 0.730 & 0.798 & Completed \\
    Qwen3 14B            & 0.649 & 0.714 & 0.789 & Completed \\
    Gemma 4 E4B           & 0.636 & 0.705 & 0.781 & Completed \\
    Llama 3.1 8B Instruct & 0.632 & 0.704 & 0.764 & Completed \\
    Gemma 4 26B A4B       & ---   & ---   & ---   & Failed \\
    \bottomrule
  \end{tabular}
\end{table}

All completed LLM runs used 4-bit quantization and LoRA with three epochs; Mistral used \texttt{max\_seq\_length} 512 in pipeline-resolved configuration while Gemma/Mistral multimodal wrappers targeted language-model projection modules.

Gemma 4 26B failed with repeated load and memory errors documented in \texttt{errors.jsonl}, including CUDA out-of-memory and meta-tensor loading failures on 2$\times$24\,GB GPUs.

\section{Stage-2 Encoder Baselines}

\subsection{BERT Token Classification}

Stage-2 trains a BERT-base-uncased token classifier as a CPU-friendly baseline (\texttt{artifacts/skillllm\_v2/bert\_baseline/metrics.json}).
Exact span F1 is 0.059 with soft-skill F1 0.019 on the SkillSpan test evaluation used for Stage-2 artifacts.

\subsection{GLiNER}

GLiNER (\texttt{urchade/gliner\_medium-v2.1}) ran in zero-shot mode because fine-tuning was skipped (\texttt{training\_mode: zero\_shot} in metrics metadata).
Exact span F1 is 0.076.

\section{Primary Generative Extractor}

The Stage-2 Mistral LoRA model was trained for 80 epochs on 4{,}800 training examples with 2$\times$ GPU distributed data parallel training (\texttt{training\_summary.json}).
Configuration specifies LoRA rank 16, alpha 32, learning rate $1.2 \times 10^{-4}$, batch size 1, gradient accumulation 32, and \texttt{max\_seq\_length} 512.

Test metrics in \texttt{artifacts/skillllm\_v2/mistral\_lora/metrics.json}:

\begin{itemize}
  \item Exact span F1: 0.670
  \item Token F1: 0.788
  \item Partial span F1: 0.814
  \item Hard-skill F1: 0.726
  \item Soft-skill F1: 0.604
\end{itemize}

These supersede the three-epoch legacy Mistral numbers for Stage-2 system claims but do not invalidate the legacy benchmark as a controlled multi-model comparison at equal epoch budget.

\section{Wrapper Variants}

Wrapper systems evaluated in Stage-2 include boundary refinement only, Mistral+GLiNER+taxonomy, and full CAG/graph pipelines.
Reported extraction metrics pass through the non-regression gate described in Chapter~\ref{ch:methodology}.

\section{Comparison Fairness}

Legacy LLM comparisons share dataset splits and evaluation scripts but differ in backbone capacity and training duration from Stage-2 80-epoch Mistral.
Encoder baselines share Stage-2 evaluation code but use different model families and task heads.
The earlier \texttt{Experiments/} resume benchmark uses different metrics and is not included in the same leaderboard.
